\section{Тема 4. Простейшие операции над массивами}

\subsection{Постановка задачи (вариант~4)}

Вычислить среднее значение $m$ и дисперсию $d$ для заданного массива
$X(n)$ наблюдений. Дисперсия вычисляется по формуле:
\[
  d = \frac{1}{n}\sum_{i=1}^{n} (x_i - m)^2.
\]

\subsection{Математическое обоснование}

Среднее арифметическое:
\[
  m = \frac{1}{n}\sum_{i=1}^{n} x_i.
\]
Дисперсия характеризует разброс значений относительно среднего.
Алгоритм требует двух проходов по массиву: первый --- для
вычисления~$m$, второй --- для суммирования квадратов отклонений.

\subsection{Блок-схема алгоритма}

\begin{center}
\begin{tikzpicture}[
  node distance=10mm, start chain=going below,
  nodes={draw=black, top color=white, bottom color=lightgray!50,
         minimum width=5.2cm, align=center, font=\small},
  arrow/.style={->, >=Stealth, thick},
]
  \node[draw, rounded corners, on chain] {Start};
  \node[shape=trapezium, trapezium left angle=70,
        trapezium right angle=110, on chain, draw, join=by arrow]
       {$n$,\ $X[0..n{-}1]$};
  \node[shape=rectangle, on chain, draw, join=by arrow]
       {$\mathrm{sum} := 0$};
  \node[shape=chamfered rectangle, chamfered rectangle xsep=10pt,
        on chain, draw, join=by arrow]
       (l1) {$i := 0,\ n{-}1$};
  \begin{scope}[local bounding box=sc1]
    \node[shape=rectangle, on chain, draw, join=by arrow]
         {$\mathrm{sum} \mathrel{+}= X[i]$};
    \coordinate[on chain, join=by arrow] (l1e);
  \end{scope}
  \draw[arrow] (sc1.south) -| ($(sc1.west)-(0.8,0)$) |- (l1);
  \coordinate[on chain] (sp1);
  \draw[arrow] (l1) -| ($(sc1.east)+(0.8,0)$) |- (sp1);
  \node[shape=rectangle, on chain, draw, join=by arrow]
       {$m := \mathrm{sum}/n$,\quad $\mathrm{dsum} := 0$};
  \node[shape=chamfered rectangle, chamfered rectangle xsep=10pt,
        on chain, draw, join=by arrow]
       (l2) {$i := 0,\ n{-}1$};
  \begin{scope}[local bounding box=sc2]
    \node[shape=rectangle, on chain, draw, join=by arrow]
         {$\mathrm{dsum} \mathrel{+}= (X[i]{-}m)^2$};
    \coordinate[on chain, join=by arrow] (l2e);
  \end{scope}
  \draw[arrow] (sc2.south) -| ($(sc2.west)-(0.8,0)$) |- (l2);
  \coordinate[on chain] (sp2);
  \draw[arrow] (l2) -| ($(sc2.east)+(0.8,0)$) |- (sp2);
  \node[shape=rectangle, on chain, draw, join=by arrow]
       {$d := \mathrm{dsum}/n$};
  \node[shape=trapezium, trapezium left angle=70,
        trapezium right angle=110, on chain, draw, join=by arrow]
       {$m,\ d$};
  \node[draw, rounded corners, on chain, join=by arrow] {End};
\end{tikzpicture}
\end{center}

\subsection{Текст программы}

\lstinputlisting[style=Cstyle]{../src/t04.c}

\subsection{Тестовые данные}

\begin{center}
\begin{tabular}{l|ll}
\toprule
Массив $X$ & $m$ & $d$ \\
\midrule
$\{1,2,3,4,5\}$ & 3.0 & 2.0 \\
$\{10,10,10,10\}$ & 10.0 & 0.0 \\
\bottomrule
\end{tabular}
\end{center}

\textbf{Результат выполнения программы} (первый тест):
\begin{lstlisting}[language={},basicstyle=\ttfamily\small,frame=single]
Введите количество элементов n: 5
Введите элементы массива:
1 2 3 4 5
Среднее значение m = 3.0000
Дисперсия d = 2.0000
\end{lstlisting}
